\section{Documentazione del Codice Sorgente}
\label{appendix:firmware}

In questa appendice viene riportata l'organizzazione modulare del firmware sviluppato per il SoC RP2040. Il codice è stato ingegnerizzato seguendo criteri di modularità e astrazione dell'hardware, per facilitare la manutenzione e il testing dei singoli sottosistemi.

\subsection{Struttura del Repository}
La gerarchia del progetto rispecchia la separazione tra configurazione dell'ambiente di build (gestita da PlatformIO), codice sorgente principale e moduli driver specifici per le periferiche.

\begin{lstlisting}[language=bash, caption={Albero del progetto firmware e descrizione dei moduli}, basicstyle=\ttfamily\scriptsize]
.
|-- Makefile             # Script per l'automazione del processo di build
|-- platformio.ini       # Configurazione ambiente, librerie e partizioni Flash
|-- README.md            # Documentazione sintetica del repository
`-- src
    |-- main.cpp         # Entry point: inizializzazione core e loop principale
    |-- minmea           # Libreria esterna per il parsing delle sentenze NMEA
    |   |-- minmea.c
    |   `-- minmea.h
    `-- Modules          # Moduli driver e logica di astrazione hardware
        |-- buttons.cpp  # Gestione ISR e debouncing asincrono dei pulsanti
        |-- buttons.h
        |-- buzzer.cpp   # Driver per segnali acustici di feedback
        |-- buzzer.h
        |-- display_ui.cpp # Logica di rendering e interfaccia grafica OLED
        |-- display_ui.h
        |-- Functions.cpp  # Funzioni helper e utility di sistema
        |-- Functions.h
        |-- gps.cpp      # Gestione UART e interfacciamento con SAM-M8Q
        |-- gps.h
        |-- leds.cpp     # Controllo striscia LED RGB (tramite offloading PIO)
        |-- leds.h
        |-- temp.cpp     # Interfaccia per sensore di temperatura ambientale
        `-- temp.h
\end{lstlisting}

\subsection{Configurazione e Riproducibilità dell'Ambiente}

Il file \texttt{platformio.ini} costituisce la \textit{Single Source of Truth} del progetto. Esso garantisce che l'intero ecosistema — comprensivo di toolchain ARM GCC, framework Arduino-Pico e dipendenze esterne — sia vincolato a versioni specifiche, assicurando la determinazione dell'ambiente di build su qualsiasi postazione di lavoro.

\subsubsection*{Gestione delle Dipendenze e del Core}
A differenza dell'ambiente Arduino IDE standard, PlatformIO permette di specificare il core di Earle Philhower (\texttt{earlephilhower}), ottimizzato per il SoC RP2040, e di definire rigorosamente il partizionamento della memoria Flash per ospitare il filesystem \textit{LittleFS}.

\begin{lstlisting}[language=C, caption={Configurazione platformio.ini per l'ottimizzazione del core RP2040.}]
[env:pico]
platform = https://github.com/maxgerhardt/platform-raspberrypi.git@1.18.0
board = nanorp2040connect
board_build.core = earlephilhower
board_build.filesystem_size = 65536 ; Riserva 64KB per LittleFS
upload_protocol = picotool

lib_deps =
    adafruit/Adafruit GFX Library @ ^1.12.4
    adafruit/Adafruit SSD1306 @ ^2.5.16
    fastled/FastLED @ ^3.10.3
\end{lstlisting}

\subsubsection*{Workflow di Sviluppo e Automazione tramite Makefile}
Sebbene PlatformIO offra una potente interfaccia a riga di comando (\textit{CLI}), è stato predisposto un \texttt{Makefile} \cite{gnu_make} per astrarre le operazioni comuni e gestire flag di compilazione condizionale. Questo permette di configurare il comportamento del firmware (es. attivazione dei log di debug o feedback acustici) senza modificare il codice sorgente.

\begin{itemize}
    \item \textbf{Compilazione Condizionale}: Tramite la variabile \texttt{PLATFORMIO\_BUILD\_FLAGS}, vengono iniettate macro nel preprocessore C++ (es. \texttt{DEBUG}, \texttt{BUZZER\_INIT\_BEEP}).
    \item \textbf{Target Specializzati}: Il comando \texttt{make debug} automatizza il ciclo build-upload-monitor, riducendo i tempi morti durante la fase di testing.
\end{itemize}

\begin{lstlisting}[language=make, caption={Makefile per la gestione dei profili di build (Silent, Debug, Production).}]
PIO = pio
DEBUG ?= 0
BEEP ?= 1

# Iniezione dei flag nel preprocessore tramite variabili d'ambiente
FLAGS = -Os -DDEBUG=$(DEBUG) -DBUZZER_INIT_BEEP=$(BEEP)
export PLATFORMIO_BUILD_FLAGS = $(FLAGS)

.PHONY: all debug fs clean

# Target standard: Compilazione e caricamento
all: 
	$(PIO) run --target upload

# Profilo Debug: Log attivi e monitoraggio seriale automatico
debug: 
	$(PIO) run --target upload
	$(PIO) device monitor

# Gestione Filesystem: Caricamento dati in LittleFS
fs:
	$(PIO) run --target uploadfs

clean:
	$(PIO) run --target clean
\end{lstlisting}

\subsubsection*{Comandi CLI Fondamentali}
Per operazioni granulari, la CLI di PlatformIO mette a disposizione i seguenti strumenti:
\begin{itemize}
    \item \texttt{pio run}: Esegue il processo di build analizzando l'albero delle dipendenze.
    \item \texttt{pio run -t buildfs}: Genera l'immagine binaria del filesystem a partire dalla cartella locale \texttt{data/}.
    \item \texttt{pio device monitor}: Avvia l'ascolto sulla porta seriale virtuale (CDC) per il debugging real-time.
\end{itemize}


\subsection{Implementazione dei Moduli Hardware}
Questa sezione raccoglie le implementazioni salienti dei driver sviluppati, con particolare enfasi sulla gestione del determinismo e delle risorse del SoC RP2040.

\subsubsection*{Sottosistema di Input: buttons.cpp}
Contiene l'implementazione delle Routine di Servizio degli Interrupt (ISR) e la logica di validazione temporale (\textit{debouncing}) tramite gli Alarm del Pico-SDK, come discusso nella Sezione 5.

% Qui inserirai \lstinputlisting{...}

\paragraph{Implementazione della logica interrupt}
Per abilitare l'interrupt sui gpio usiamo la funzione \textit{gpio\_set\_irq\_enabled\_with\_callback} del pico-sdk (\href{https://www.raspberrypi.com/documentation/pico-sdk/hardware.html}{hardware/gpio.h})

\begin{lstlisting}[language=C]
pinMode(SWpin[i], INPUT_PULLUP);

gpio_set_irq_enabled_with_callback(SWpin[i], GPIO_IRQ_LEVEL_HIGH, false, &buttons_callback);
gpio_set_irq_enabled_with_callback(SWpin[i], GPIO_IRQ_LEVEL_LOW, false, &buttons_callback);
gpio_set_irq_enabled_with_callback(SWpin[i], GPIO_IRQ_EDGE_RISE, false, &buttons_callback);
gpio_set_irq_enabled_with_callback(SWpin[i], GPIO_IRQ_EDGE_FALL, true, &buttons_callback);
\end{lstlisting}

\textbf{NOTA}: come misura di sicurezza vengono disabilitati gli altri trigger che non ci interessano, infatti potenzialmente possono essercene diversi attivi simultaneamente.\\

Essendo i segnali dei pulsanti attivi bassi, il trigger dell'interrupt avviene per il Falling Edge.
La logica utilizzata per l'handling degli interrupts é la seguente:
una volta rilevato il fronte di discesa (\textit{falling edge}), il processore sospende l'esecuzione del task corrente per saltare alla routine di servizio (\textit{ISR}) denominata \texttt{buttons\_callback}. Poiché l'architettura dell'RP2040 raggruppa gli eventi GPIO in un unico segnale verso il controller NVIC (\textit{Nested Vectored Interrupt Controller}), la funzione riceve come parametri l'identificativo del pin (\texttt{gpio}) e l'evento che ha scatenato l'interrupt (\texttt{events}).
\begin{lstlisting}[language=C]
void buttons_callback(uint gpio, uint32_t events) {
    if(events == GPIO_IRQ_EDGE_RISE)
    {
        gpio_set_irq_enabled(gpio, GPIO_IRQ_EDGE_RISE, false);
        gpio_set_irq_enabled(gpio, GPIO_IRQ_EDGE_FALL, true);
        return;
    }
    switch(gpio){

        case SW_UP:
        case SW_DOWN:
        case SW_LEFT:
        case SW_RIGHT:
        {
            gpio_set_irq_enabled(gpio, GPIO_IRQ_EDGE_FALL, false);
            add_alarm_in_ms (DEBOUNCE_DELAY, &debounce_callback, (void*) gpio, true);
        } break; 
        default:
        {
            ;
        }
    } return; 
}
\end{lstlisting}

All'interno dell'handler, viene implementata una struttura di controllo (switch case) per discriminare la natura dell'evento; per ora si prega di ignorare il primo costrutto if, in quanto la sua utilità sarà piú chiara successivamente. 
Lo switch case non discrimina il singolo pulsante fra i quattro, ma si occupa di attivare un alarm interrupt, ossia un delay non bloccante per il processore che attiva una ISR (debounce\_callback) dopo un certo delay (=DEBOUNCE\_DELAY). É importante notare che viene anche disabilitato l'interrupt sul falling edge, per evitare che un bounce del segnale possa generare molti interrupts. 

Analizziamo ora l'handling dell'alarm interrupt:

\begin{lstlisting}[language=C]
long long debounce_callback(alarm_id_t id, void *user_data){
    uint PinNum=(uint) user_data;

    if(!gpio_get( PinNum ))
    {
        switch( PinNum )
        {
            case SW_UP: up_press=true; break;

            case SW_DOWN: down_press=true; break;
            
            case SW_LEFT: left_press=true; break;

            case SW_RIGHT: right_press=true; break;

            default:
            {
                break;
            }
        }  
            gpio_set_irq_enabled(PinNum, GPIO_IRQ_EDGE_RISE, true);
            return 0;
    }
    gpio_set_irq_enabled(PinNum, GPIO_IRQ_EDGE_FALL, true); 
    return 0;
}
\end{lstlisting}

In questo caso viene verificato (tramite \textit{gpio\_get(PinNum)}) se il Pin dopo il delay \textit{DEBOUNCE\_DELAY} é ancora basso:
\begin{itemize}
    \item se ancora basso allora si determina il gpio dell'interrupt (variabile \textit{user data} passata dalla funzione \textit{buttons\_callback}), e scrive nella variabile bool true.
    Ora viene attivato l'interrupt per il fronte di salita, per identificare il rilascio del pulsante. L'handler di tale fronte é ancora \textit{buttons\_callback}.
    
    \item se il pulsante é stato rilasciato in un tempo inferiore a \textit{DEBOUNCE\_DELAY}, allora viene riattivato l'interrupt per il fronte di discesa senza fare nient'altro.
\end{itemize}

Ritornando al primo if-case della funzione \textit{buttons\_callback}:
\begin{lstlisting}[language=C]
void buttons_callback(uint gpio, uint32_t events) {
    if(events == GPIO_IRQ_EDGE_RISE)
    {
        gpio_set_irq_enabled(gpio, GPIO_IRQ_EDGE_RISE, false);
        gpio_set_irq_enabled(gpio, GPIO_IRQ_EDGE_FALL, true);
        return;
    }
    ...
\end{lstlisting}

Serve per identificare il primo fronte di salita (e quindi il rilascio) dopo una pressione \textbf{valida} dei pulsanti; infatti viene disattivato l'interrupt per il fronte di salita e attivato quello per il fronte di discesa, senza fare nient'altro. 

Questa scelta implementativa é stata presa per evitare che qualche rumore dato dal rilascio del pulsante possa attivare un'interrupt che causerebbe solo una falsa pressione, nel caso sfortunato in cui il rumore perdurasse per un tempo \textit{DEBOUNCE\_DELAY}.

Il main loop verifica in maniera sincrona la pressione dei pulsanti tramite le variabili bool che ne indicano la avvenuta pressione. In questo modo separiamo le decisioni software dovute dalla pressione dei pulsanti (esecuzione lenta) da quella di verifica della pressione degli stessi (esecuzione rapida).

Queste scelte mantengono il \textit{Core} libero da cicli di \textit{polling} intensivi, delegando la logica di debouncing e l'attesa di un evento direttamente all'hardware del SoC. Tuttavia é necessaria dell'attenzione aggiuntiva per gestire situazioni in cui  gli interrupt generati potrebbero interrompere fasi critiche del programma, ad esempio il flusso di dati proveniente dal modulo GNSS SAM-M8Q o lo scambio dati con il PC.




\subsubsection*{Interfaccia GNSS: gps.cpp}
Implementa la ricezione seriale asincrona e il parsing deterministico delle sentenze NMEA (\texttt{\$GPRMC}, \texttt{\$GPGGA}) tramite la libreria \textit{minmea}.

% Qui inserirai \lstinputlisting{...}


PARSING STRINGHE NMEA
\begin{lstlisting}[language=C]
// Esempio di parsing di una sentenza RMC
if (minmea_parse_rmc(&frame, line)) {
    float lat = minmea_tocoord(&frame.latitude);
    float lon = minmea_tocoord(&frame.longitude);
    float speed = minmea_tofloat(&frame.speed);
}
\end{lstlisting}

\subsubsection*{Schermo}
\begin{lstlisting}[language=C]
#include <Adafruit_SSD1306.h>
Adafruit_SSD1306 display(SCREEN_WIDTH, SCREEN_HEIGHT, &Wire, -1);

// Inizializzazione con indirizzo I2C
if(!display.begin(SSD1306_SWITCHCAPVCC, 0x3C)) { 
    Serial.println(F("Errore allocazione SSD1306"));
}
display.clearDisplay(); // Svuota il buffer interno
\end{lstlisting}


\subsubsection*{LED}
\begin{lstlisting}[language=C]
#include <FastLED.h>
CRGB leds[NUM_LEDS];

// Impostazione colore tramite spazio HSV (Tonalita, Saturazione, Valore)
leds[0] = CHSV(hue, 255, 255);
FastLED.show();
\end{lstlisting}

\subsubsection*{File System}

Per importare la libreria é sufficiente
\begin{lstlisting}[language=C]
#include "LittleFS.h" 
\end{lstlisting}
e per montare il filesystem:
\begin{lstlisting}[language=C]
LittleFS.begin();
\end{lstlisting}
Di default se il mount fallisce la libreria formatta la parte di flash riservata al filesystem e riprova il mount.

Altre funzioni utilizzate sono:

\begin{lstlisting}[language=C]
//Apre un file in maniera simile 
//ad fopen (ritorna stream pointer)
File f = LittleFS.open(path, mode);

//Rimuove un file 
//(ritorna bool)
LittleFS.remove(path);

//Verifica se il path esiste 
//(ritorna bool)
LittleFS.exists(path);

//Crea una cartella 
//(ritorna bool)
LittleFS.mkdir(path);

//Rimuove una cartella 
//(ritorna bool)
LittleFS.rmdir(path);
\end{lstlisting}






CAPIRE DOVE METERE:
Il firmware interagisce con il PIO tramite una semplice interfaccia di scrittura:

\begin{lstlisting}[language=C]
// Caricamento del dato nel buffer FIFO della State Machine
pio_sm_put_blocking(pio0, sm, color_data << 8u);
\end{lstlisting}